\section*{Abstract}

This manuscript studies finite-time breakdown claims for the three-dimensional
incompressible Navier--Stokes equations from smooth divergence-free data.  The
object under test is always the same object: the classical solution generated by
the initial data up to the alleged terminal time, together with the terminal
record offered as the reason that same solution cannot be continued.  A terminal
claim is not carried by a point alone.  It must be carried by positive-radius
velocity fields around that point as the radii shrink.

The paper keeps two certification layers separate.  The direct layer tries to
prove estimates strong enough to continue the original solution.  In the current
proof state, that route reaches a specific annular obstruction: finite physical
energy can still allow source mass to concentrate on shrinking heat-time windows,
leaving a scale-critical terminal residue.  The paper states the direct theorem
still needed to rule out that residue, instead of treating the unfinished
estimate as a proof.

The contrapositive layer is the class-membership test.  Once a surviving record
\(Q\) is selected and admitted as the record of the same Navier--Stokes branch,
it has to carry a positive same-fluid carrier, participation in the same
pressure-viscosity-incompressibility law, and one coherent positive-scale
velocity field.  These are the Pack, Part, and Field faces.  Failure of the
first required face gives a class exit,
\[
  \Exit(Q;\mathfrak O_{\NS}^{\mathrm{work}})=
  \neg\Member(Q;\mathfrak O_{\NS}^{\mathrm{work}}),
\]
and the object cannot serve as an in-class nonsmooth terminal continuation of
the same solution.

The manuscript is organized by branch family.  Each branch states its direct
aim, carries out the relevant calculation, identifies the stopping obstruction,
records the survivor object, and then reads that object through the
class-membership test.  The periodic zero-force problem on \(\mathbb T^3\) is
the first concrete domain treated here.  Whole-space \(\mathbb R^3\) export is
tracked as its own branch and scope boundary.

\section{Scope And Proof Polarity}
\label{sec:scope-proof-polarity}

The Navier--Stokes global regularity problem is usually approached in the direct direction.  That route tries to estimate the original solution strongly enough that finite-time loss of smoothness has no room to occur.  That is the gold route in this paper.  It is the most direct certification layer because it keeps the solution itself inside the continuation argument.

The present proof also uses a contrapositive layer.  It keeps the same equation and the same branch.  It begins with the direct attempts because those attempts identify the actual terminal objects that threaten continuation.  When the forward accounting leaves a precise obstruction, the proof asks whether that obstruction is still a same-solution Navier--Stokes terminal record.  If it is admitted to the CM test and cannot remain on the member-continuation side, the first failed Pack, Part, or Field face yields class exit.

This polarity matters for reading the branch material.  A failed gold estimate becomes useful when it leaves a selected terminal record tied to the original branch.  A CM face failure has a different job: it is the reason the alleged finite nonsmooth terminal record cannot also be the same in-class Navier--Stokes member record.

The first concrete domain is the periodic zero-force problem on \(\mathbb T^3\).  Whole-space \(\mathbb R^3\) material is tracked as its own export and exterior-source boundary.  Comparison material and Euler boundary material are included only where they supply a same-solution proof role or a smoothness-boundary role.  A branch-local result reaches the full Clay conclusion only through the export argument that carries it there.
